\section{速度から位置と道のりへ}\label{節:運動と積分}

一直線上の運動では，位置を $s(t)$，速度を $v(t)=s'(t)$，加速度を $a(t)=v'(t)$ とする。速度が連続なら，微積分の基本定理により
\[
s(t_1)-s(t_0)=\int_{t_0}^{t_1}v(t)\dd t.
\]
左辺は位置の差，すなわち変位である。出発点に戻れば零になるが，途中で移動した距離が零という意味ではない。

速度が正なら正の向きに進み，負なら負の向きに進む。移動した距離を向きによらず加えるには速さ $|v(t)|$ を積分して，
\begin{equation}\label{式:道のり}
\text{道のり}=\int_{t_0}^{t_1}|v(t)|\dd t
\end{equation}
とする。速さは非負，速度には向きに対応する符号がある。この違いは，面積と符号付きの定積分の違いと同じである。

速度が零になる時刻は，進む向きが変わる候補である。ただし零になるだけでは反転したとは限らない。たとえば $v(t)=(t-1)^2$ は $t=1$ で零になるが，前後で非負である。道のりを計算するときは，零点で区間を分けたうえで，各区間の符号を確かめる。

\begin{example}{変位と道のりを分けて求める}
時刻 $0\leq t\leq4$ において，速度が $v(t)=(t-1)(t-3)$ で与えられる。初期位置は $s(0)=0$ とする。$t=4$ の位置，道のり，平均速度，平均の速さを求めよ。
\end{example}

速度を積分して初期条件を使うと，
\[
s(t)=\frac{t^3}{3}-2t^2+3t.
\]
したがって $s(4)=\frac43$ である。道のりでは速度の符号を調べる。$v$ は $0\leq t<1$ で正，$1<t<3$ で負，$3<t\leq4$ で正である。折り返す位置は
\[
s(1)=\frac43,\qquad s(3)=0
\]
なので，道のりは
\[
\left|\frac43-0\right|+\left|0-\frac43\right|
 +\left|\frac43-0\right|=4.
\]
これは $\int_0^1v\dd t-\int_1^3v\dd t+\int_3^4v\dd t$ と計算しても同じである。全時間は $4$ だから，平均速度は $\cfrac{\cfrac43}{4}=\frac13$，平均の速さは $\frac44=1$ となる。

加速度は $a(t)=2t-4$ である。加速度が正でも，いつも速さが増えるわけではない。$v\ne0$ の時刻では
\[
\frac{\dd}{\dd t}|v(t)|=
\begin{cases}
a(t),&v(t)>0,\\
-a(t),&v(t)<0
\end{cases}
\]
だから，速度と加速度が同符号なら速さが増え，異符号なら減る。この例では $1<t<2$ と $3<t<4$ で速さが増える。負の向きに進みながら速くなることもあるのである。

平面上の運動 $\bigl(x(t),y(t)\bigr)$ では，速度は二成分 $\bigl(x'(t),y'(t)\bigr)$ を持ち，速さは $\sqrt{x'(t)^2+y'(t)^2}$ となる。その積分が道のりになることは，式\eqref{式:媒介曲線の弧長}そのものである。運動では同じ場所を再び通った分も道のりに含むので，曲線を重複してたどることは誤りではなく，数えるべき内容になる。

